% ============================== CHAPTER 2 ===================================
\chapter{Literature Survey}
\label{ch:literature}

This chapter surveys the research that \setu{} builds on and reacts against.
The survey has three purposes. First, it establishes what is technically
available for compressing a translation model to edge size, for training a
compact student from a strong teacher, and for covering many Indian languages at
once. Second, it compares those available technologies so that the choices made
in Chapter~\ref{ch:design} can be justified rather than asserted. Third, it
isolates the specific capability that none of the surveyed systems provides:
offline, on-device, multi-directional Indic translation within 200\,MB, which
becomes the problem statement of this project.

Twenty papers published from 2019 onward are reviewed, ordered by year of
publication. They fall into five groups: compact and mobile neural machine
translation, knowledge distillation, preference optimisation and alignment,
multilingual and Indic-specific translation, and offline or edge deployment.

\section{Introduction to the Survey}
\label{sec:litintro}

Neural machine translation converged on the Transformer architecture, and with
it on a scaling regime in which quality is bought with parameters and data.
That regime produced IndicTrans2, which translates all twenty-two scheduled
Indian languages at high quality but requires GPU-class hardware. The research
directions that matter for this project are the ones that ask what can be
recovered after the parameters are taken away again.

Three lines of work answer that question. Compression research (pruning,
quantisation, shallow architectures) reduces a trained model to a size a
device can hold. Distillation research transfers behaviour from a large teacher
to a small student, either by matching output distributions or by training the
student on the teacher's decoded output. Alignment research, developed for large
language models, replaces a single reference target with a preference signal
over several candidate outputs. The central question that this project settles
empirically is which of the latter two gives a very small translation student
the better training signal when the available parallel data is web-mined and
noisy.

A fourth line, multilingual translation, matters for coverage. Shared
multilingual models exploit transfer between related languages, which is
valuable for the Indian language family, but they also introduce interference
between directions and pivoting error when no direct model exists.

\section{Summary of Papers Surveyed}
\label{sec:papers}

\begin{longtable}{@{}p{0.4cm} p{3.6cm} p{2.1cm} p{1.5cm} p{5.6cm}@{}}
\caption{Summary of the literature surveyed (2019--2025), in ascending order of year}
\label{tab:litsurvey}\\
\toprule
\rowcolor{headblue}
\textbf{\#} & \textbf{Title} & \textbf{Authors} & \textbf{Venue \& Year} & \textbf{Contributions and drawbacks} \\
\midrule
\endfirsthead
\multicolumn{5}{c}{\textit{Table \thetable{} (continued)}}\\
\toprule
\rowcolor{headblue}
\textbf{\#} & \textbf{Title} & \textbf{Authors} & \textbf{Venue \& Year} & \textbf{Contributions and drawbacks} \\
\midrule
\endhead
\midrule
\multicolumn{5}{r}{\textit{continued on next page}}\\
\endfoot
\bottomrule
\endlastfoot

1 & DistilBERT, a distilled version of BERT: smaller, faster, cheaper and lighter
  & V. Sanh \textit{et al.}
  & arXiv, 2019
  & Halves the layer count of BERT via distillation while retaining about 97\,\% of
    language-understanding performance, and establishes triple-loss distillation
    as a practical recipe. \textbf{Drawback:} demonstrated for encoders on
    classification, not for autoregressive sequence generation, and a measurable
    accuracy loss remains on complex tasks. \\

2 & Neural Machine Translation on Mobile Devices
  & B. Li \textit{et al.}
  & IEEE, 2020
  & Compresses NMT models for mobile and edge deployment, cutting memory
    footprint and inference latency to make real-time on-device translation
    feasible. \textbf{Drawback:} the reduction techniques degrade translation
    quality, and the degradation is most visible on long and syntactically
    complex sentences. \\

3 & SMaLL-100: Introducing Shallow Multilingual Machine Translation for
    Low-Resource Languages
  & A. Mohammadshahi \textit{et al.}
  & Springer, 2020
  & A compact multilingual translator covering many low-resource directions with
    substantially lower memory use and faster inference than M2M-100.
    \textbf{Drawback:} quality still trails the large model it compresses,
    particularly on high-resource pairs where capacity matters. \\

4 & The Tatoeba Translation Challenge -- Realistic Data Sets for Low Resource
    and Multilingual MT
  & J. Tiedemann
  & WMT, 2021
  & Releases a benchmark spanning 555 languages, enabling systematic evaluation
    of low-resource, zero-shot and few-shot translation on realistic data.
    \textbf{Drawback:} the data is noisy and inconsistent across languages, and
    exploiting it fully still requires large-capacity models. \\

5 & Shallow-to-Deep Training for Neural Machine Translation
  & B. Li \textit{et al.}
  & ACL, 2021
  & Introduces incremental shallow-to-deep training, stacking Transformer layers
    progressively to improve convergence speed and training stability at depth.
    \textbf{Drawback:} the final models remain computationally heavy and the
    staged schedule adds engineering complexity. \\

6 & Nearest Neighbor Knowledge Distillation for Neural Machine Translation
  & Z. Yang, R. Sun, X. Wan
  & IEEE, 2022
  & Uses $k$-nearest-neighbour retrieval over the training set to supply
    additional supervision, improving quality in low-resource conditions.
    \textbf{Drawback:} the datastore is memory-intensive and retrieval slows both
    training and inference, which is disqualifying for edge deployment. \\

7 & Speech-to-Speech Translation Improvements with Self-Supervised Learning
  & C. Xu \textit{et al.}
  & IEEE, 2022
  & Combines self-supervised pre-training, multitask learning and synthetic data
    to obtain large BLEU gains in speech-to-speech translation, especially for
    low-resource pairs. \textbf{Drawback:} depends on large-scale pre-training and
    a multi-stage pipeline that is impractical to deploy on a device. \\

8 & Direct Speech Translation
  & Y. Jia \textit{et al.}
  & Springer, 2023
  & Surveys end-to-end speech translation and characterises its principal
    obstacles: data scarcity, latency and bias. \textbf{Drawback:} a survey rather
    than a system; no concrete model or experimental validation is offered. \\

9 & Direct Preference Optimization: Your Language Model is Secretly a Reward Model
  & R. Rafailov \textit{et al.}
  & NeurIPS, 2023
  & Reformulates preference alignment as a simple classification loss over
    preferred and dispreferred responses, removing the separate reward model and
    the reinforcement-learning loop of RLHF, with better stability and lower
    cost. \textbf{Drawback:} effectiveness depends heavily on the quality and
    diversity of the preference data; the paper studies alignment of large
    language models, not distillation into small ones. \\

10 & All-in-One Knowledge Distillation Framework for Neural Machine Translation
  & Z. Miao \textit{et al.}
  & EMNLP, 2023
  & Distils into several student models simultaneously within one framework,
    improving training efficiency and enabling scalable multilingual
    distillation. \textbf{Drawback:} joint optimisation across students creates
    gradient conflicts, and low-resource directions are the ones that lose. \\

11 & Continual Knowledge Distillation for Neural Machine Translation
  & Y. Zhang, P. Li, M. Sun, Y. Liu
  & ACL, 2023
  & Adds new languages to an existing multilingual model incrementally without
    forgetting previously learned ones, supporting long-lived multilingual
    deployments. \textbf{Drawback:} training complexity grows with the number of
    tasks and the stability-plasticity balance requires careful tuning. \\

12 & IndicTrans2: Towards High-Quality and Accessible Machine Translation Models
     for All 22 Scheduled Indian Languages
  & J. Gala \textit{et al.}
  & TMLR, 2023
  & Delivers open translation models covering all twenty-two scheduled Indian
    languages, together with the BPCC corpus and the IN22 benchmark; the current
    reference point for Indic MT. \textbf{Drawback:} the flagship models need
    GPU-class compute and multiple gigabytes of memory, which rules out offline
    edge deployment. \\

13 & Multilingual Indian Language Neural Machine Translation System Using mT5
     Transformer
  & A. Jha \textit{et al.}
  & IEEE, 2023
  & Fine-tunes mT5 for Indian-language translation and reports strong BLEU on the
    directions covered, confirming that large pre-trained multilingual encoders
    transfer to Indic data. \textbf{Drawback:} evaluated on only a few languages
    and built on a large pre-trained backbone, so it scales neither to all 22
    languages nor to on-device inference. \\

14 & IWSLT 2024 Low-Resource Speech Translation Submissions
  & M. Zafar \textit{et al.}
  & IWSLT, 2024
  & Shows that cascaded ASR-then-MT pipelines still outperform end-to-end speech
    translation in genuinely low-resource settings, and reports competitive BLEU.
    \textbf{Drawback:} cascading compounds recognition errors into translation
    errors; end-to-end alternatives remain data-starved. \\

15 & Direct Preference Optimization for Neural Machine Translation with Minimum
     Bayes Risk Decoding
  & G. Yang, J. Chen, W. Lin, B. Byrne
  & NAACL, 2024
  & Applies DPO to NMT using MBR decoding to construct preference pairs, showing
    that preference learning can improve translation quality with relatively
    little data. \textbf{Drawback:} evaluation does not extend to extremely
    low-resource languages, and the models studied are far larger than an
    edge-deployable student. \\

16 & Multilingual NMT System for English to Low-Resource Indic Languages
     (Assamese and Bengali)
  & K. Kashyap, S. K. Sarma
  & IEEE, 2024
  & A one-to-many multilingual system that improves low-resource quality through
    shared representations and back-translated synthetic data.
    \textbf{Drawback:} synthetic data introduces its own noise, and coverage is
    limited to two languages in one direction. \\

17 & Multilingual Neural Machine Translation System for Indic-to-Indic Languages
  & S. B. Das \textit{et al.}
  & Science\-Direct, 2024
  & Translates directly among several Indic languages using shared models and
    pivoting, addressing a direction pair largely ignored by English-centric
    systems. \textbf{Drawback:} error propagates through the pivot language and
    accuracy degrades relative to direct models. \\

18 & SURYA-TAC: LLM-Based Offline Portable Speech-to-Speech Translation System
  & P. S. Bohra, K. Berwal
  & IEEE, 2024
  & Proposes an offline, edge-deployed tactical translation system for use in
    harsh environments without connectivity, the closest published motivation
    to this project. \textbf{Drawback:} largely conceptual, without real-world
    validation, and constrained by data scarcity for the target languages. \\

19 & Latency-Aware Simultaneous Machine Translation Using Large Language Models
  & B. Fu \textit{et al.}
  & Springer, 2024
  & Introduces latency-aware token generation so that large language models can
    perform simultaneous translation while maintaining accuracy.
    \textbf{Drawback:} computationally expensive and premised on models far too
    large for a handset. \\

20 & Offline English-to-Hindi Statistical Machine Translation System
  & S. Gangwar, I. Naim
  & Science\-Direct, 2024
  & Implements genuinely offline English-to-Hindi translation using the Moses SMT
    toolkit, demonstrating that offline operation is achievable at low cost.
    \textbf{Drawback:} statistical MT is substantially weaker than neural
    approaches, and the system covers a single direction of a single language
    pair. \\

\end{longtable}

\subsection{Discussion of Key Papers}
\label{subsec:keypapers}

Four of the surveyed papers determined the shape of this project and merit
closer discussion.

\textbf{IndicTrans2} \cite{ref:indictrans2} is the teacher. Its contribution is
not only the models but the surrounding artifacts: the BPCC corpus and the IN22
benchmark, which together made Indic MT reproducible. \setu{} adopts its
script-qualified FLORES language codes (\texttt{hin\_Deva}, \texttt{eng\_Latn})
as internal identifiers so that student tokenizer tags remain stable as
languages are added. The project's quality target (80\,\% of teacher BLEU) is
defined relative to this system. A practical constraint arose here: the official
\texttt{ai4bharat} checkpoints and the BPCC dataset are access-gated on Hugging
Face, returning HTTP 401 without an authenticated account that has accepted the
terms. The project therefore uses the ungated rotary distilled-200M checkpoints
released by IndicTrans2 co-author Raj Dabre, and the ungated Samanantar corpus.
This lowers the achievable ceiling, and Chapter~\ref{ch:results} quantifies by
how much.

\textbf{Direct Preference Optimization} \cite{ref:dpo} supplied the project's
original hypothesis. Its key insight (that the optimal policy under a
Bradley-Terry preference model has a closed form, so preference learning reduces
to a classification loss on preferred and dispreferred pairs) makes preference
training cheap enough to attempt on a student model. \setu{} implements DPO with
a frozen supervised-fine-tuning reference model, chrF-ranked teacher hypotheses
as preferred candidates, and $k$-nearest-neighbour retrieved translations as
additional negatives. The implementation works as specified: the margin accuracy
reaches 1.0, meaning the trained policy ranks every preferred candidate above
its dispreferred counterpart. It nonetheless fails to improve translation
quality over sequence-level distillation, for reasons developed in
Section~\ref{sec:proposedsolution} and Chapter~\ref{ch:results}.

\textbf{Sequence-level knowledge distillation}, introduced by Kim and Rush
\cite{ref:seqkd} and instantiated in this project's baseline, replaces the human
reference target with the teacher's own one-best decoded output. It is the
approach that won. The mechanism is best understood as denoising: the teacher's
output is internally consistent and drawn from a single distribution, whereas
web-mined references are inconsistent in register, alignment quality and
fidelity. A 52-million-parameter student has no spare capacity to average over
that inconsistency.

\textbf{DistilBERT} \cite{ref:distilbert} and \textbf{SMaLL-100}
\cite{ref:small100} between them establish the size-quality frontier that this
project operates on. DistilBERT shows that roughly 97\,\% of a large model's
performance survives halving its depth; SMaLL-100 shows the same for
multilingual translation specifically. Both suggest that a compact student can
retain most of a teacher's ability, which is what the 0.80 ratio target
formalises.

\section{Comparative Study of Available Technologies}
\label{sec:comparative}

Table~\ref{tab:techcompare} compares the candidate technologies for each design
decision the project faced, and records the choice made together with its
justification.

\begin{longtable}{@{}p{2.5cm} p{3.8cm} p{3.85cm} p{3.8cm}@{}}
\caption{Comparative study of available technologies and the choices made}
\label{tab:techcompare}\\
\toprule
\rowcolor{headblue}
\textbf{Decision} & \textbf{Alternatives considered} & \textbf{Trade-off} & \textbf{Choice and justification} \\
\midrule
\endfirsthead
\multicolumn{4}{c}{\textit{Table \thetable{} (continued)}}\\
\toprule
\rowcolor{headblue}
\textbf{Decision} & \textbf{Alternatives considered} & \textbf{Trade-off} & \textbf{Choice and justification} \\
\midrule
\endhead
\midrule
\multicolumn{4}{r}{\textit{continued on next page}}\\
\endfoot
\bottomrule
\endlastfoot

Teacher model
& IndicTrans2 1B; IndicTrans2 distilled-200M; mT5-based Indic NMT; NLLB-200
& The 1B flagship gives the highest ceiling but is gated and needs roughly 5\,GB
  for weights alone; mT5 covers fewer Indic languages
& \textbf{Ungated rotary IndicTrans2 distilled-200M.} Accessible without an
  authenticated account, fits available memory, and covers all 22 scheduled
  languages. The lower ceiling is reported rather than hidden \\

Training signal for the student
& Human references (plain SFT); teacher one-best (\seqkd{}); chrF-ranked
  preference pairs with DPO; \seqkd{} followed by DPO
& References are abundant but noisy; teacher output is clean but caps the
  student at teacher quality; preference pairs add ranking information but need
  a reliable ranker
& \textbf{\seqkd{} on teacher one-best}, selected by the controlled experiment in
  Section~\ref{sec:proposedsolution}: 17.09 BLEU against 11.10 for
  reference-trained DPO at matched size \\

Student architecture
& Marian encoder--decoder; distilled Transformer; shallow multilingual
  (SMaLL-100 style); decoder-only LM
& Encoder--decoder remains the strongest inductive bias for translation;
  decoder-only needs far more parameters for equivalent quality
& \textbf{Marian-style Transformer, 6 encoder and 6 decoder layers, $d=512$,
  52.6\,M parameters.} Quantises to 103.9\,MB, comfortably inside the 200\,MB
  budget \\

Vocabulary
& Shared 32{,}000-token SentencePiece; 16{,}000-token SentencePiece; byte-level BPE
& A larger vocabulary shortens sequences but leaves rare tokens undertrained on
  a corpus of this size
& \textbf{16{,}000-token SentencePiece.} A 32{,}000-token vocabulary was measured to
  produce degenerate decoding on this data volume (Section~\ref{sec:runlog}) \\

Multilingual strategy
& One model per direction; one shared multilingual model; shared model with
  All-in-One KD
& Per-direction models avoid interference and permit independent quality gates;
  a shared model saves storage but suffers low-resource interference
  \cite{ref:aiokd}
& \textbf{One student per direction}, with the All-in-One KD mechanism
  implemented and tested for future consolidation \\

Indic-to-Indic coverage
& Train $n(n-1)$ direct models; pivot through English; use a gated
  Indic-to-Indic teacher
& Direct models scale quadratically at 104\,MB each; pivoting doubles latency
  and can compound errors; the Indic-to-Indic teacher is gated
& \textbf{English pivot.} $2n$ models serve $n^2$ directions; measured two-hop
  latency remains around 50--65\,ms \\

Language expansion
& Full retraining; All-in-One KD; Elastic Weight Consolidation
  \cite{ref:continualkd}
& Retraining is prohibitive; EWC protects old languages via a Fisher-weighted
  penalty
& \textbf{EWC implemented and validated} as the incremental path; per-direction
  training used for Phase~1 \\

Quantisation and runtime
& PyTorch on-device; TensorFlow Lite; ONNX Runtime with INT8/INT4; GGML
& ONNX Runtime has mature CPU kernels, a small footprint and wide platform
  support; its CPU provider lacks a true 4-bit kernel
& \textbf{ONNX Runtime with progressive INT8 then INT4 export via Optimum.}
  411\,MB FP32 compresses to 103.9\,MB \\

Evaluation metric
& BLEU; chrF/chrF++; COMET; human evaluation
& BLEU is standard but penalises agglutinative languages; COMET correlates
  better with human judgement but needs a neural model at evaluation time
& \textbf{BLEU and chrF via sacreBLEU}, reported as a student/teacher ratio;
  COMET implemented for future benchmarking \\

\end{longtable}

\subsection{Why These Technologies Were Selected}
\label{subsec:why}

The selection logic follows from the four objectives in
Section~\ref{sec:objectives} treated as hard constraints rather than
aspirations.

The 200\,MB size budget eliminates every alternative that keeps a teacher-sized
model on the device, and therefore mandates distillation into a student of
roughly 50 million parameters. That parameter budget in turn eliminates
retrieval-augmented approaches such as nearest-neighbour distillation
\cite{ref:knnkd}, whose datastore alone would exceed the budget.

The 500\,ms latency budget on CPU eliminates beam search at deployment (the
deployed engine decodes greedily, with beam search reserved for offline
evaluation) and eliminates any method requiring a retrieval lookup per token.

The offline requirement eliminates any runtime that fetches weights, tokenizer
artifacts or telemetry, which is why the tokenizer is a local SentencePiece
model rather than a Hugging Face identifier resolved at load time.

The quality target, finally, is what forced the empirical comparison. Both
sequence-level distillation and preference distillation were plausible routes to
80\,\% of teacher BLEU, and the literature did not settle the question for
students this small on data this noisy. The project therefore built both and
measured them.

\section{Drawbacks of the Existing System}
\label{sec:drawbacks}

The survey exposes five drawbacks that, taken together, leave the offline Indic
translation problem unsolved.

\textbf{Cloud dependency.} Almost every high-quality MT system in production
requires an active network connection. IndicTrans2 \cite{ref:indictrans2}, mT5
based systems \cite{ref:mt5indic} and LLM-based simultaneous translators
\cite{ref:latency} all assume server-side execution. In regions with intermittent
connectivity, which is where the language barrier is most acute, these
systems are unavailable exactly when needed. Cloud execution also means that
every translated sentence is transmitted to and processed by a third party.

\textbf{High compute requirements.} The models that achieve state-of-the-art
Indic quality hold hundreds of millions to billions of parameters and need
GPU-class memory. The IndicTrans2 flagship requires roughly 5\,GB for weights
alone. Mid-range handsets and ARM single-board computers cannot host them, and
the compression literature \cite{ref:mobilenmt, ref:small100} consistently reports
quality degradation as the price of shrinking them.

\textbf{Error propagation in cascaded pipelines.} Systems that chain automatic
speech recognition to translation \cite{ref:iwslt} compound recognition errors
into translation errors. The IWSLT 2024 results show cascades still outperforming
end-to-end models in low-resource settings, which means the compounding problem
is currently unavoidable in that design.

\textbf{Narrow language coverage.} Compact and specialised systems cover very few
languages. The mT5-based Indic system \cite{ref:mt5indic} evaluates on a handful;
the Assamese-Bengali system \cite{ref:assamese} covers two in one direction; the
offline SMT system \cite{ref:offlinesmt} covers a single direction of a single
pair. Systems with broad coverage are precisely the ones too large to deploy.

\textbf{Data scarcity for low-resource languages.} Languages such as Kashmiri,
Dogri, Konkani, Bodo and Santali lack parallel corpora at the scale neural
translation requires. Where corpora do exist they are frequently web-mined and
noisy, as the Tatoeba challenge \cite{ref:tatoeba} documents. This project
encountered the problem directly: only eleven of the twenty-two scheduled
languages had ungated data sufficient to train against.

\section{Problem Statement}
\label{sec:problem}

\begin{quote}
\itshape
There is a need for a lightweight, offline-capable multilingual translation
framework that delivers high-quality translation across the scheduled languages
of India on edge devices (retaining at least 80\,\% of a strong teacher's BLEU
within a 200\,MB artifact and under 500\,ms of CPU latency per sentence, with
zero network calls at inference) and that supports incremental addition of new
languages without catastrophic forgetting and without a quadratic growth in the
number of trained models.
\end{quote}

The overall problem decomposes into six sub-problems, each addressed by an
identifiable part of the system.

\textbf{P1: Corpus quality.} The available ungated parallel data (Samanantar)
is web-mined, and contains misaligned pairs, length-ratio anomalies,
wrong-script segments and duplicates. \emph{Sub-problem:} construct a cleaning
pipeline that yields a corpus a 52-million-parameter student can learn from.
Addressed by the Corpus Loader module (Section~\ref{subsec:mod1}).

\textbf{P2: Choice of training signal.} It is not established which signal
(human references, teacher one-best output, or preference-ranked teacher
hypotheses) best trains a student an order of magnitude smaller than its
teacher on noisy data. \emph{Sub-problem:} construct all three and compare them
under matched conditions. Addressed by the Teacher Distiller and Trainer modules
and reported in Section~\ref{sec:headtohead}.

\textbf{P3: Training stability.} A Transformer trained from scratch at this
scale is unstable: three distinct failure modes (empty output, unconditional
decoding, and mid-training collapse) were encountered. \emph{Sub-problem:}
identify and eliminate each. Addressed in Section~\ref{sec:runlog}.

\textbf{P4: Deployment footprint.} A 52.6-million-parameter FP32 model
occupies 411\,MB, over twice the size budget, and decodes too slowly in PyTorch
on CPU. \emph{Sub-problem:} compress and accelerate without materially losing
quality. Addressed by the Quantiser module (Section~\ref{subsec:mod4}).

\textbf{P5: Indic-to-Indic coverage.} No ungated Indic-to-Indic teacher exists,
and training direct models for every pair would require $n(n-1)$ models at
104\,MB each. \emph{Sub-problem:} serve every pair without quadratic growth.
Addressed by the English-pivot router (Section~\ref{subsec:mod5}).

\textbf{P6: Extensibility.} Adding a language must not require retraining
everything or degrade the languages already supported. \emph{Sub-problem:}
provide a mechanism for incremental expansion with measured retention. Addressed
by the All-in-One KD and Elastic Weight Consolidation implementations
(Section~\ref{subsec:mod8}).

\section{Proposed Solution}
\label{sec:proposedsolution}

\setu{} proposes a four-stage pipeline (distil, train, quantise, serve)
instantiated once per translation direction and bound together by a routing
layer that pivots through English.

\textbf{Stage 1: Teacher distillation.} The ungated IndicTrans2 distilled-200M
teacher decodes every source sentence in the cleaned corpus, and its one-best
output replaces the human reference as the training target. This is
sequence-level knowledge distillation in the sense of Kim and Rush
\cite{ref:seqkd}. Greedy decoding is used, which is four to five times faster than
beam search and produces targets of equivalent utility.

\textbf{Stage 2: Student training.} A 52.6-million-parameter Marian
encoder--decoder with a 16{,}000-token SentencePiece vocabulary is trained from
scratch on the distilled corpus with token-level cross-entropy, padding masked
out of the loss, linear learning-rate warmup and decay, label smoothing of 0.1,
dropout of 0.2 and gradient-norm clipping at 1.0. Every one of these components
was added in response to a measured failure, not adopted by default;
Section~\ref{sec:runlog} gives the history.

\textbf{Stage 3: Quantisation.} The trained student is exported to ONNX,
validated against ONNX Runtime, and quantised progressively to INT8 and then
INT4, with latency and size benchmarked at each stage. The result is a 103.9\,MB
artifact, a fourfold compression from 411\,MB.

\textbf{Stage 4: Offline serving.} An inference engine loads the quantised ONNX
model and a local SentencePiece tokenizer and decodes greedily. A router
resolves each request to a directly trained model where one exists, and
otherwise chains two students through English. One engine serves a REST API, a
command-line tool, an offline progressive web application and a browser client.

\subsection{The Deciding Experiment}
\label{subsec:deciding}

The project's founding hypothesis was H1: \emph{preference distillation produces
better compact students than sequence-level knowledge distillation}. It was
tested at matched student size, on identical data, with evaluation on real human
references held out from training, a design that is fair regardless of what
each system trained on. Table~\ref{tab:h1} gives the outcome.

\begin{table}[H]
\centering
\caption{Head-to-head at 100{,}000 training sentences, Hindi\arrowto{}English,
identical 52.6\,M student, evaluated on held-out human references}
\label{tab:h1}
\begin{tabular}{@{}l l c c c@{}}
\toprule
\rowcolor{headblue}
\textbf{System} & \textbf{Training target} & \textbf{BLEU} & \textbf{Ratio} & \textbf{Verdict} \\
\midrule
S0: SFT on references & human references (Samanantar) & 10.95 & 0.389 & baseline \\
\rowcolor{rowgrey}
S1: \seqkd{}          & teacher one-best              & \textbf{17.09} & \textbf{0.607} & \textbf{best} \\
S2: DPO-distill       & references, then DPO          & 11.10 & 0.394 & no gain over S0 \\
\bottomrule
\end{tabular}
\end{table}

H1 is refuted. Sequence-level distillation beats preference distillation by
6.0 BLEU, an improvement of about 54\,\%, and DPO applied on top of a
reference-trained base adds 0.15 BLEU over that base, indistinguishable from
noise. The replacement hypothesis, H1$'$, is supported by every subsequent
result: \emph{reference noise, not the training objective, is what caps a
compact Indic student.} Teacher-distilled targets remove that noise, and the
student then continues to improve with data: the ratio rises from 0.607 at
100{,}000 sentences to 0.764 at 250{,}000 and 0.806 at 500{,}000, whereas
reference-trained students plateau near 0.66 regardless of data volume.

The DPO implementation is not defective, which is what makes the result
informative. Its margin accuracy reaches 1.0, meaning it ranks preferred above
dispreferred candidates perfectly; it consistently improves chrF by one to two
points while costing about one BLEU point. The preference signal is being
learned correctly; it simply is not the signal that a 52-million-parameter
student most needs when its reference targets are noisy. \setu{} therefore ships
\seqkd{}, and reports DPO as a controlled negative result.
